% Generated draft: requires mathematical and source review.
\documentclass[11pt]{article}
\usepackage[margin=25mm]{geometry}
\usepackage{fontspec}
\setmainfont{Latin Modern Roman}
\usepackage{amsmath,amssymb,mathtools,unicode-math}
\setmathfont{Latin Modern Math}
\usepackage{enumitem,needspace,xurl,hyperref}
\setlength{\parindent}{0pt}
\setlength{\parskip}{5pt}
\setlength{\emergencystretch}{3em}
\newcommand{\N}{\mathbb{N}}
\newcommand{\Z}{\mathbb{Z}}
\newcommand{\Q}{\mathbb{Q}}
\newcommand{\R}{\mathbb{R}}
\newcommand{\C}{\mathbb{C}}
\newcommand{\F}{\mathbb{F}}
\newcommand{\A}{\mathbb{A}}
\newcommand{\PP}{\mathbb{P}}
\newcommand{\E}{\mathbb{E}}
\newcommand{\eps}{\varepsilon}
\newcommand{\dd}{\,\mathrm{d}}
\begin{document}
\section*{UnsolvedMath ID: 30002838}
\textbf{Source problem code:} catalogue entry
\subsection{Definitions and mathematical statement}
Let \(\mathrm{TMF}\) denote the spectrum named in the supplied record, and let \(\pi_n(\mathrm{TMF})\) denote its degree-\(n\) homotopy group. The precise version of this spectrum and the intended range of degrees are not specified. For a nonnegative integer \(n\), let
\[
\mathcal L_n=\{\mathfrak g:\mathfrak g\text{ is a real semisimple Lie algebra of compact type},\ \dim_{\R}\mathfrak g=n\}/\cong.
\]
Here compact type means that \(\mathfrak g\) is the Lie algebra of a compact Lie group, and \(\cong\) denotes Lie algebra isomorphism. Using isomorphism classes is an explicit interpretation of the proposed parameters; the source does not specify the parameter space or the meaning of naturality.

For an \(n\)-dimensional compact Lie group \(G\), an ordered basis \((v_1,\ldots,v_n)\) of \(T_eG\) defines a left-invariant tangent framing by
\[
\lambda_g=\bigl((\dd L_g)_e v_1,\ldots,(\dd L_g)_e v_n\bigr),\qquad g\in G,
\]
where \(L_g\) is left multiplication. Under the usual interpretation of a framed manifold as a stable homotopy class, and provided the intended \(\mathrm{TMF}\) has a specified unit \(\eta:\mathbb S\to\mathrm{TMF}\), write
\[
\mathcal C_n=\left\{\eta_*[G,\lambda]\in\pi_n(\mathrm{TMF}):G\text{ is a compact Lie group of dimension }n,\ \lambda\text{ is a left-invariant framing}\right\}.
\]
Here \(\mathbb S\) is the sphere spectrum and \([G,\lambda]\in\pi_n(\mathbb S)\) is the framed bordism class. This makes explicit a conventional interpretation of the phrase ``classes represented by compact Lie groups with their left-invariant framings''; the record does not specify this map, restrict connectedness, or fix framing choices.

\textbf{Question and conjectural target.} Does the proposed Stolz--Teichner cocycle construction produce a natural family of elements
\[
\Theta_n:\mathcal L_n\longrightarrow\pi_n(\mathrm{TMF})?
\]
The original wording asks for such a family and describes its realization by the cocycle construction as plausible. It does not explicitly require \(\Theta_n\) to be injective, surjective, or a bijection onto all of \(\pi_n(\mathrm{TMF})\).

The further supplied conjecture is that the relevant construction hits exactly \(\mathcal C_n\). If its degree-\(n\) output is denoted schematically by
\[
F_n:\mathcal D_n\longrightarrow\pi_n(\mathrm{TMF}),
\]
this target is
\[
\operatorname{im}(F_n)=\mathcal C_n.
\]
Neither \(\mathcal D_n\) nor the precise cocycle construction is supplied. In particular, the record does not determine whether \(F_n\) means only the semisimple Lie algebra family \(\Theta_n\) or a broader construction. Thus the displayed equality is a conditional formulation of the stated target, not a fully specified conjecture or an asserted result.
\subsection{Short English statement}
The problem asks whether \(n\)-dimensional semisimple compact-type Lie algebras naturally give classes in \(\pi_n(\mathrm{TMF})\) through a proposed Stolz--Teichner cocycle construction, and conjectures that the relevant construction produces exactly the classes represented by compact Lie groups with left-invariant framings. The construction, its domain, and the strength of the proposed parametrization remain unspecified. The supplied open status is dated curation, not an independently verified conclusion.
\subsection{Sources}
\begin{itemize}
\item Local registry: \texttt{data/unsolvedmath/problems.json}, permanent UnsolvedMath ID \(30002838\), source problem code \texttt{catalogue entry}, title ``Lie-Algebraic Parametrization of TMF Classes.'' The supplied record assigns proposed year \(2015\) and supplies no named proposer.
\item Supplied report citation: \emph{Subfactors and Conformal Field Theory} (2016), \emph{Oberwolfach Reports}, pp.\ 849--926, DOI \texttt{10.4171/owr/2015/16}. \url{https://doi.org/10.4171/owr/2015/16}. The background describes these reports as records of mathematical workshops and their open-problem discussions.
\item Supplied literature source, labeled ``Henriques--Douglas, proposed construction'': \url{https://doi.org/10.4171/owr/2015/16}.
\item Supplied literature source, labeled ``Douglas--Henriques, TMF via conformal nets'': \url{https://arxiv.org/abs/1103.4187}.
\item Literature assessment checked at \texttt{2026-08-21}: the supplied curation labels the problem open, describes compact Lie algebras as yielding a proposed family of Stolz--Teichner cocycles, and describes the exact-image assertion for framed compact Lie groups as conjectural. This is supplied dated curation derived from the OWR record, not an independently verified assessment. The supplied background expressly requests independent verification of time-sensitive status and cited claims.
\item The registry labels the statement \texttt{corrected\_verified} and reports a manual semantic, typographic, and LaTeX review claiming meaning-preserving cleanup. These are supplied editorial claims; they do not establish mathematical correctness or resolve the formulation issues recorded here.
\end{itemize}
\end{document}
